% ============================================================
%  Objetivo 4 — De humanos a bacterias y viceversa
% ============================================================

La proteína humana que buscamos como homóloga es la \textbf{guanilato ciclasa soluble (sGC)},
cuyas subunidades \textbf{GUCY1A1} y \textbf{GUCY1B1} contienen el dominio HNOB en el
extremo N-terminal.

% ============================================================
\section{Ejercicio 4a --- Búsqueda de homólogos humanos por distintos métodos}

\begin{consigna}
Buscar homólogos de la Proteína X en humanos usando: a) BLAST, b) PSI-BLAST, c) Pfam,
d) PROSITE. ¿Cuál es el método más sensible? ¿Con cuál llegó más rápido al resultado?
\end{consigna}

\subsection{a) BLAST}

Corrimos BLASTp de la secuencia de \textit{C. tetani} contra nr restringido a
\textit{Homo sapiens} con parámetros por defecto. La sGC \textbf{no apareció} entre los
hits significativos. La identidad entre el dominio HNOB de \textit{C. tetani} y el de la
sGC humana es $\sim$20--25\%, que está por debajo del umbral de sensibilidad de BLAST
estándar. Con tan baja identidad, los HSPs no alcanzan scores suficientemente altos como
para ser significativos estadísticamente.

\subsection{b) PSI-BLAST}

Corrimos PSI-BLAST con la secuencia completa de 617 aa contra nr restringido a \textit{Homo
sapiens}. En la \textbf{iteración 1} no hubo hits significativos (solo 2 lamininas por debajo
del umbral). A partir de la \textbf{iteración 2} empezaron a aparecer muchos hits, pero todos
correspondían a proteínas coiled-coil humanas: lamininas, miosinas, nineínas, plectinas y
desmoplaquinas. En las iteraciones 3 y 4 la PSSM quedó completamente dominada por esas
proteínas, y la sGC (\textbf{GUCY1A1/GUCY1B1}) no apareció en ninguna iteración.

El problema es que la Proteína X tiene dos dominios: el HNOB (posiciones 1--190), que es el
que nos conectaría con la sGC, y el MCPsignal (posiciones 331--617), que tiene estructura de
coiled-coil. Como PSI-BLAST buscó con la secuencia completa, los hits del dominio MCPsignal
dominaron la PSSM desde la iteración 2, arrastrando la búsqueda hacia proteínas estructuralmente
similares al MCPsignal y alejándola del dominio HNOB. Este fenómeno se llama \textit{PSSM drift}:
la matriz se ``contamina'' con información de falsos positivos y la búsqueda diverge del
homólogo de interés.

Para haber encontrado la sGC con PSI-BLAST habría que buscar solo con la región HNOB (posiciones
1--190 de la secuencia), de modo que la PSSM se construya únicamente con información del dominio
sensor de NO y no se vea contaminada por los coiled-coils del MCPsignal.

\subsection{c) Pfam / InterPro}

Este fue el método más directo: al buscar la secuencia en InterPro el dominio HNOB
(PF07730) queda identificado de inmediato, y desde la página del dominio en Pfam se puede
navegar directamente a las proteínas humanas con ese dominio en la sección \textit{Taxonomy}.
Ahí aparecen las subunidades de la sGC (GUCY1A1 y GUCY1B1) sin necesidad de correr ninguna
búsqueda adicional.

\subsection{d) PROSITE}

El perfil PS50111 que detectó PROSITE en el Objetivo 3 corresponde al dominio MCPsignal,
no al HNOB, así que no recupera la sGC humana directamente. Para llegar a ella por PROSITE
habría que buscar el perfil específico de HNOB. Es más sensible que BLAST pero menos práctico
que Pfam para este caso concreto.

\subsection{Comparación de métodos}

\begin{table}[H]
\centering
\footnotesize
\begin{tabular}{lP{3cm}P{2.5cm}P{3.5cm}}
\toprule
Método & ¿Encuentra sGC? & Sensibilidad & Velocidad \\
\midrule
BLAST         & No                    & Baja para homólogos distantes & Rápido \\
PSI-BLAST     & No (PSSM drift)       & Alta (si se usa solo HNOB)    & Moderado \\
Pfam/InterPro & Sí, inmediatamente    & Muy alta                      & Muy rápido (navegación directa) \\
PROSITE       & No con PS50111        & Alta (con el perfil HNOB)     & Moderado \\
\bottomrule
\end{tabular}
\caption{Comparación de métodos para detectar la sGC humana como homóloga de la Proteína X.}
\end{table}

El \textbf{método más sensible y más rápido} para detectar la sGC resultó ser \textbf{Pfam/InterPro}:
navegando directamente por el dominio HNOB (PF07730) aparecen las subunidades de la sGC sin
necesidad de correr ninguna búsqueda iterativa. PSI-BLAST en principio es más sensible que BLAST
estándar, pero con la secuencia completa de 617 aa sufrió \textit{PSSM drift}: el dominio
MCPsignal (coiled-coil) dominó la matriz desde la iteración 2 y la sGC no apareció en ninguna de
las 4 iteraciones que corrimos. Para que PSI-BLAST funcione en este caso habría que usarlo
únicamente con la región HNOB (posiciones 1--190). BLAST estándar también falla porque la
identidad entre la Proteína X y la sGC ($\sim$20\%) está por debajo de su umbral de sensibilidad.

% ============================================================
\section{Ejercicio 4b --- Función propuesta y experimento}

\begin{consigna}
¿Qué función le asignaría a la Proteína X? Proponga un experimento para confirmarla.
\end{consigna}

A partir de todo lo analizado, la Proteína X de \textit{Clostridium tetani} sería una
\textbf{hemoproteína sensora de óxido nítrico acoplada al sistema de quimiotaxis bacteriana}.
Tiene un dominio HNOB N-terminal que une un grupo hemo y detecta NO en el entorno, y un dominio
MCPsignal C-terminal que transduciría esa señal al sistema de quimiotaxis. La función propuesta es
que cuando el sistema inmune del huésped produce NO durante la infección, la bacteria lo detecta a
través de esta proteína y regula su movimiento en respuesta. Sería una función análoga a la de la
sGC en mamíferos, que también es activada por NO vía su dominio HNOB, aunque el efector posterior
es distinto (ciclasa en animales, dominio MCP en bacterias).

Para confirmar esta función proponemos dos experimentos:

\begin{enumerate}
  \item \textbf{Reconstitución con hemo y espectroscopía UV-Vis.} Expresar el dominio HNOB
        (posiciones 1--190) de forma recombinante en \textit{E. coli}, reconstituir con hemo
        exógeno y medir el espectro de absorción (300--700\,nm) en presencia y ausencia de NO.
        Cuando el NO se une al Fe(II) del hemo se produce un Soret shift característico
        ($\sim$431 $\rightarrow$ 399\,nm). Si se observa ese desplazamiento, queda demostrado que
        la proteína une NO a través de su grupo hemo, tal como se predice del análisis de dominios.

  \item \textbf{Complementación funcional en \textit{E. coli}.} Usar una cepa de \textit{E. coli}
        con deleción de todos sus receptores MCP nativos (\textit{$\Delta$tsr $\Delta$tar $\Delta$trg
        $\Delta$tap $\Delta$aer}) y expresar en ella la Proteína X de \textit{C. tetani}. Medir la
        quimiotaxis en placas de motilidad o en ensayos de capilares con gradientes de NO. Si la
        bacteria recupera capacidad de quimiotaxis dirigida hacia o desde fuentes de NO, quedaría
        confirmado que el dominio HNOB actúa como sensor y el MCPsignal como efector funcional.
\end{enumerate}
